\begin{helper}
Let \(\eta,\varepsilon_1,\varepsilon_2,\beta\in\mathbb R\) and
\(n_0\in\mathbb N\).  If \(\beta=1/2\) and
\(\noderef{ParameterHierarchy}(\eta,\varepsilon_1,\varepsilon_2,n_0)\), then
the \(\eta\)-scale regularity event
\[
\noderef{TwoBiteRegularityEvent}
\bigl(k=\noderef{TwoBiteNaturalIndependenceScale}(1+\eta,n),
\eta,\varepsilon_1,\varepsilon_2,\omega\bigr)
\]
holds with high probability under the two-bite law
\[
\noderef{TwoBiteEventProbability}
\left(n,\noderef{TwoBiteNaturalReducedVertexCount}(n),
\noderef{TwoBiteEdgeProbability}(\beta,n),-\right).
\]
\end{helper}
\begin{proof}
Write
\[
M(n)=\noderef{TwoBiteNaturalReducedVertexCount}(n),\qquad
K(n)=\noderef{TwoBiteNaturalIndependenceScale}(1+\eta,n),
\]
and
\[
\mathbb P_n(E)=
\noderef{TwoBiteEventProbability}
\bigl(n,M(n),\noderef{TwoBiteEdgeProbability}(\beta,n),E\bigr).
\]
Unfolding \(\noderef{TwoBiteRegularityEvent}\), the target regularity event is
the intersection of the three events
\[
D_n(\omega)=\noderef{FiberAndDegreeEvent}(\omega,\varepsilon_2),
\]
\[
M_n(\omega)=
\noderef{TwoBiteMediumClosedPairsEvent}(k=K(n),\eta,\varepsilon_1,\omega),
\]
and
\[
S_n(\omega)=
\noderef{TwoBiteSmallClosedPairsEvent}(k=K(n),\eta,\varepsilon_1,\omega).
\]

The hierarchy hypothesis begins with \(0<\varepsilon_2\).  Hence
\(\noderef{FiberAndDegreeConcentration}\), with the function
\(n\mapsto M(n)\) and the equality \(\beta=1/2\), gives
\(\mathbb P_n(D_n)\to1\).

Apply \(\noderef{MediumClosedPairsBound}\) with parameters
\(\eta,\varepsilon_1,\varepsilon_2,\beta,n_0\), the same function
\(n\mapsto M(n)\), and the function \(n\mapsto K(n)\).  Its two function
identities are tautological from the definitions of \(M\) and \(K\), and its
hierarchy and \(\beta=1/2\) assumptions are exactly the standing hypotheses.
Therefore \(\mathbb P_n(M_n)\to1\).  The same argument with
\(\noderef{SmallClosedPairsBound}\) gives \(\mathbb P_n(S_n)\to1\).

It remains only to pass from the three high-probability events to their
intersection.  For any tolerance \(\rho>0\), choose thresholds after which
each of \(D_n,M_n,S_n\) has probability greater than \(1-\rho/3\).  The
complement of \(D_n\cap M_n\cap S_n\) is contained in
\(D_n^c\cup M_n^c\cup S_n^c\).  Since
\(\noderef{TwoBiteEventProbability}\) is a finite sum of nonnegative sample
weights, the finite union bound gives
\[
\mathbb P_n((D_n\cap M_n\cap S_n)^c)
\le
\mathbb P_n(D_n^c)+\mathbb P_n(M_n^c)+\mathbb P_n(S_n^c)<\rho.
\]
Equivalently, \(\mathbb P_n(D_n\cap M_n\cap S_n)>1-\rho\) eventually.  This
is precisely convergence of the regularity-event probability to \(1\).
\end{proof}
